\section{Discusión prevista y amenazas a la validez}
\subsection{Validez interna}
La amenaza principal es comparar plataformas con correcciones distintas.
Se controla derivando variantes de una base común y registrando sus cambios.
La segunda es debilitar el baseline: B3 debe reutilizar \(S_a\), disponer de
MUL y usar una disposición de datos comparable. El despacho de zero-points,
colas y epílogo tampoco puede desaparecer de la región medida.

La base docente requiere ampliar verificación. Un resultado numérico final
correcto no detecta necesariamente efectos duplicados, instrucciones ilegales
ignoradas o contadores incorrectos. El plan incluye controles de retiro,
selección de fuentes y detección de instrucciones fuera del subconjunto.
La medición se habilita solo después de cerrar esos riesgos para los kernels.

\subsection{Validez de las métricas y alcance externo}
Menos instrucciones no implica menos ciclos; menos ciclos no implica menor
tiempo si cambia la frecuencia; menos tiempo tampoco prueba menor energía.
Se evita usar esas magnitudes como sinónimos. El tiempo que tarda el simulador
en el computador anfitrión no es el tiempo del procesador modelado.

Productos punto sintéticos y matrices pequeñas permiten aislar causas, pero
no representan por sí solos una aplicación de IA ni un LLM. Tampoco una mejora
en memoria combinacional predice automáticamente el resultado con BRAM o
cache. Las afirmaciones se limitarán al dominio efectivamente evaluado.

Con \(z=8\), la recodificación a S4 puede ofrecer otra ruta eficiente. Con
muchas filas, la amortización de \(S_a\) puede favorecer a B3. Con metadatos
dinámicos, el despacho puede consumir los ahorros de D. Estos casos son
posibles explicaciones que habrá que contrastar, no conclusiones anticipadas.

\subsection{Criterios de discusión de los resultados}
Para cada diferencia se buscará un mecanismo: extracción, instrucciones de
corrección, carga de operandos, reutilización, stalls, camino crítico o código
especializado. Se contrastarán desgloses con el total del kernel. Si no aparece
ventaja, se documentará qué costo la elimina y hasta dónde puede generalizarse.
No se cambiará a posteriori la pregunta para ocultar una hipótesis refutada.

\subsection{Integridad del artefacto y publicación}
La exportación histórica no contiene un archivo LICENSE en su raíz. La
publicación del derivado requiere resolver permisos, contribuciones y licencia
con las personas que participaron en el proyecto de curso. Autoría, afiliación
y venue no se asignan en esta versión. Una plantilla editorial específica
se adoptará cuando exista una convocatoria apropiada y evidencia suficiente.
